% Ad Familiares 11.18 --- headnote
% ad-familiares-11-18, 1 June 43 BC, Rome

\textit{Cicero to D. Brutus, \textit{imperator} and consul-
designate, from Rome on 1 June 43 BC --- Perseus dateline
\textit{Scr. Romae x iiii K. Iun. a. 711 (43)}. The dateline
is corrupt in the manuscript (literally ``14 days before
the Kalends of June'' = 19 May, which is impossible given
the contents); the standard reading, following Shackleton
Bailey, takes the closing \textit{x iiii K. Iun.} as
\textit{prid. K. Iun.}, the eve of the Kalends, and dates
the letter to 31 May or 1 June 43 BC. The reply is to a set
of instructions Brutus had sent to the Senate by his
legates Galba and Volumnius, urging caution.}

\textit{Cicero is at his most public-Roman here. The
Senate, he tells Brutus, will not be patronised: it judges
him the bravest man of any who ever lived, and resents
being read in return as faint of heart. But by the time
this letter reaches camp the ground will already have
shifted under it. Lepidus's army has gone over to Antony
at the bridge over the Argenteus on 30 May; the
\textit{senatus consultum ultimum} declaring Lepidus a
public enemy will follow on 30 June. Cicero's rhetorical
question --- ``who could think him so frenzied that, in
the most longed-for peace, he would declare war on the
state?'' --- is being answered in Narbonese Gaul as he
writes. The third paragraph's closing tricolon --- that
the state lacks neither counsel, nor courage, nor a
commander while Brutus lives --- is the public note
Cicero will keep sounding through June, even as the
private letters grow sharper.}
